\documentclass[border=3pt]{standalone}
\usepackage{luatexja-fontspec}
\setmainjfont[Scale=1,BoldFont=HiraMinProN-W6]{Hiragino Mincho ProN}
\usepackage{newtxtext}
\usepackage{newtxmath}
\ltjsetparameter{xkanjiskip=0pt}
\usepackage{booktabs,xcolor}
\definecolor{tred}{RGB}{192,57,43}
\definecolor{tnavy}{RGB}{31,58,122}
\newcommand{\R}[1]{\textcolor{tred}{#1}}
\newcommand{\B}[1]{\textcolor{tnavy}{\textbf{#1}}}
\newcommand{\U}[1]{\textcolor{tnavy}{\underline{#1}}}
\newcommand{\dg}{\rlap{$^{\dagger}$}}
\renewcommand{\arraystretch}{1.2}
\setlength{\tabcolsep}{5pt}
\begin{document}%
\begin{tabular}{lll}%
\toprule
分割 & 画像 & 用途\\
\midrule
train & 分布内の学習画像5万枚（ImageNet-64は1万枚） & 比較手法の当てはめ\\
val1 & 分布内の画像2,000枚（ImageNet-64は1,500枚） & OLRの$\vec{\mu}$、$\mathbf{S}$、$\vec{\mu}_{o}$。比較手法の当てはめ\\
val2 & 分布内と各分布外集合から1,000枚ずつ（FashionMNISTは2,000枚） & OLRの開発\\
test & 500枚ずつ & 比較手法との比較\\
test2〜6 & 500枚ずつ（test6のCIFAR-10は258枚） & 分割を変えた確認\\
\bottomrule
\end{tabular}%
\end{document}
